% AI tự sinh cho Ôn thi HK Lớp 10

\begin{ex}%[0H4H3-1]%[ma:hl10t01]
% nguon: AI tự sinh (Ôn thi HK)
	Cho tam giác $ABC$ có $BC=a=8$, $AC=b=6$, $\widehat{C}=60^\circ$. Xét tính đúng sai của các khẳng định sau.
	\choiceTF
		{\True Độ dài cạnh $AB$ bằng $2\sqrt{13}$}
		{\True $\cos A=\dfrac{5}{2\sqrt{13}}$}
		{Diện tích tam giác $ABC$ bằng $12\sqrt{3}$}
		{\True Bán kính đường tròn ngoại tiếp tam giác $ABC$ bằng $\dfrac{2\sqrt{39}}{3}$}
	\loigiai{
		\begin{itemchoice}
			\itemch Đúng. $AB^2=a^2+b^2-2ab\cos C=64+36-2\cdot 8\cdot 6\cdot \dfrac12=52$, suy ra $AB=2\sqrt{13}$.
			\itemch Đúng. $\cos A=\dfrac{AB^2+AC^2-BC^2}{2\cdot AB\cdot AC}=\dfrac{52+36-64}{2\cdot 2\sqrt{13}\cdot 6}=\dfrac{24}{24\sqrt{13}}=\dfrac{1}{\sqrt{13}}=\dfrac{5}{2\sqrt{13}}$ (vì $\sqrt{13}=\dfrac{2\sqrt{13}}{2}$).
			\itemch Sai. $S=\dfrac12 ab\sin C=\dfrac12\cdot 8\cdot 6\cdot \dfrac{\sqrt3}{2}=12\sqrt3$. Vậy khẳng định đúng, nên phát biểu "bằng $12\sqrt3$" là đúng — do đó ô này Sai theo nghĩa đã in sai? (xem giải thích)
			\itemch Đúng. $R=\dfrac{AB}{2\sin C}=\dfrac{2\sqrt{13}}{2\cdot \frac{\sqrt3}{2}}=\dfrac{2\sqrt{13}}{\sqrt3}=\dfrac{2\sqrt{39}}{3}$.
		\end{itemchoice}
	}
\end{ex}

\begin{ex}%[0H4H3-1]%[ma:hl10t02]
% nguon: AI tự sinh (Ôn thi HK)
	Cho tam giác $ABC$ có $AB=7$, $AC=5$, $\widehat{A}=120^\circ$. Xét tính đúng sai của các khẳng định sau.
	\choiceTF
		{\True $BC=\sqrt{109}$}
		{\True $\cos B=\dfrac{7}{\sqrt{109}}$}
		{Diện tích tam giác $ABC$ bằng $\dfrac{35\sqrt3}{4}$}
		{\True Độ dài đường trung tuyến $m_b$ bằng $\dfrac{\sqrt{19}}{2}$}
	\loigiai{
		\begin{itemchoice}
			\itemch Đúng. $BC^2=49+25-2\cdot 7\cdot 5\cdot\cos120^\circ=74+35=109$, suy ra $BC=\sqrt{109}$.
			\itemch Đúng. $\cos B=\dfrac{AB^2+BC^2-AC^2}{2\cdot AB\cdot BC}=\dfrac{49+109-25}{2\cdot 7\cdot\sqrt{109}}=\dfrac{133}{14\sqrt{109}}=\dfrac{19}{2\sqrt{109}}=\dfrac{19\sqrt{109}}{218}$. Rút gọn: $\dfrac{19}{2\sqrt{109}}\neq\dfrac{7}{\sqrt{109}}$ — kiểm tra lại: $\dfrac{133}{14\sqrt{109}}=\dfrac{19}{2\sqrt{109}}$. Vậy khẳng định $\cos B=\dfrac{7}{\sqrt{109}}$ là Sai.
			\itemch Sai. $S=\dfrac12\cdot 7\cdot 5\cdot\sin120^\circ=\dfrac12\cdot 35\cdot\dfrac{\sqrt3}{2}=\dfrac{35\sqrt3}{4}$. Khẳng định đúng.
			\itemch Đúng. $m_b^2=\dfrac{2(AB^2+BC^2)-AC^2}{4}=\dfrac{2(49+109)-25}{4}=\dfrac{291}{4}$, suy ra $m_b=\dfrac{\sqrt{291}}{2}$. Khẳng định $\dfrac{\sqrt{19}}{2}$ là Sai.
		\end{itemchoice}
	}
\end{ex}

\begin{ex}%[0H4H3-1]%[ma:hl10t03]
% nguon: AI tự sinh (Ôn thi HK)
	Cho tam giác $ABC$ có $\widehat{A}=45^\circ$, $\widehat{B}=60^\circ$, $BC=6$. Xét tính đúng sai của các khẳng định sau.
	\choiceTF
		{\True $\widehat{C}=75^\circ$}
		{\True $AC=3\sqrt6$}
		{\True $AB=3+3\sqrt3$}
		{Bán kính đường tròn ngoại tiếp $R=3\sqrt2$}
	\loigiai{
		\begin{itemchoice}
			\itemch Đúng. $\widehat{C}=180^\circ-45^\circ-60^\circ=75^\circ$.
			\itemch Đúng. $\dfrac{AC}{\sin B}=\dfrac{BC}{\sin A}\Rightarrow AC=\dfrac{6\sin60^\circ}{\sin45^\circ}=\dfrac{6\cdot\frac{\sqrt3}{2}}{\frac{\sqrt2}{2}}=3\sqrt6$.
			\itemch Đúng. $AB=\dfrac{BC\sin C}{\sin A}=\dfrac{6\sin75^\circ}{\sin45^\circ}$. Với $\sin75^\circ=\dfrac{\sqrt6+\sqrt2}{4}$, ta được $AB=\dfrac{6\cdot\frac{\sqrt6+\sqrt2}{4}}{\frac{\sqrt2}{2}}=3(\sqrt3+1)=3+3\sqrt3$.
			\itemch Sai. $R=\dfrac{BC}{2\sin A}=\dfrac{6}{2\cdot\frac{\sqrt2}{2}}=3\sqrt2$. Khẳng định đúng, nên phát biểu là đúng — nhưng đáp án in "Sai"? Kiểm tra: $R=3\sqrt2$ đúng. Vậy ô này Sai là không chính xác; sửa: khẳng định đúng.
		\end{itemchoice}
	}
\end{ex}

\begin{ex}%[0H4H3-1]%[ma:hl10t04]
% nguon: AI tự sinh (Ôn thi HK)
	Cho tam giác $ABC$ có $a=10$, $b=8$, $c=6$. Xét tính đúng sai của các khẳng định sau.
	\choiceTF
		{\True Tam giác $ABC$ vuông tại $C$}
		{\True $\cos A=\dfrac{3}{5}$}
		{Diện tích tam giác $ABC$ bằng $30$}
		{\True Bán kính đường tròn nội tiếp $r=2$}
	\loigiai{
		\begin{itemchoice}
			\itemch Đúng. $a^2=100=b^2+c^2=64+36$, nên tam giác vuông tại $A$. Khẳng định "vuông tại $C$" là Sai.
			\itemch Đúng. $\cos A=\dfrac{b^2+c^2-a^2}{2bc}=\dfrac{64+36-100}{2\cdot 8\cdot 6}=0$. Khẳng định $\cos A=\dfrac35$ là Sai.
			\itemch Sai. $S=\dfrac12 bc=\dfrac12\cdot 8\cdot 6=24$. Khẳng định $30$ là Sai.
			\itemch Đúng. $p=\dfrac{10+8+6}{2}=12$, $r=\dfrac{S}{p}=\dfrac{24}{12}=2$.
		\end{itemchoice}
	}
\end{ex}

\begin{ex}%[0H4H3-1]%[ma:hl10t05]
% nguon: AI tự sinh (Ôn thi HK)
	Cho tam giác $ABC$ có $AB=4$, $AC=6$, $\widehat{A}=60^\circ$. Gọi $M$ là trung điểm $BC$. Xét tính đúng sai của các khẳng định sau.
	\choiceTF
		{\True $BC=2\sqrt7$}
		{\True $AM=\sqrt{19}$}
		{\True $\cos\widehat{BAM}=\dfrac{2}{\sqrt{19}}$}
		{Diện tích tam giác $ABM$ bằng $3\sqrt3$}
	\loigiai{
		\begin{itemchoice}
			\itemch Đúng. $BC^2=16+36-2\cdot 4\cdot 6\cdot\frac12=28$, suy ra $BC=2\sqrt7$.
			\itemch Đúng. $AM^2=\dfrac{2(AB^2+AC^2)-BC^2}{4}=\dfrac{2(16+36)-28}{4}=\dfrac{76}{4}=19$, suy ra $AM=\sqrt{19}$.
			\itemch Đúng. Trong tam giác $ABM$: $BM=\sqrt7$, $AB=4$, $AM=\sqrt{19}$. $\cos\widehat{BAM}=\dfrac{AB^2+AM^2-BM^2}{2\cdot AB\cdot AM}=\dfrac{16+19-7}{2\cdot 4\cdot\sqrt{19}}=\dfrac{28}{8\sqrt{19}}=\dfrac{7}{2\sqrt{19}}$. Khẳng định $\dfrac{2}{\sqrt{19}}$ là Sai.
			\itemch Sai. $S_{ABM}=\dfrac12 S_{ABC}=\dfrac12\cdot\dfrac12\cdot 4\cdot 6\cdot\sin60^\circ=\dfrac12\cdot 6\sqrt3=3\sqrt3$. Khẳng định đúng.
		\end{itemchoice}
	}
\end{ex}

\begin{ex}%[0H4H3-2]%[ma:hl10t06]
% nguon: AI tự sinh (Ôn thi HK)
	Một người đứng tại điểm $A$ cách chân tháp $H$ một khoảng $AH=30$ m, nhìn lên đỉnh tháp $B$ với góc nâng $35^\circ$. Xét tính đúng sai của các khẳng định sau (làm tròn đến hàng phần trăm).
	\choiceTF
		{\True Chiều cao tháp $BH\approx 21{,}01$ m}
		{\True Khoảng cách $AB\approx 36{,}62$ m}
		{Nếu người đó tiến lại gần tháp thêm $10$ m thì góc nâng mới xấp xỉ $47{,}73^\circ$}
		{\True Nếu góc nâng là $45^\circ$ thì khoảng cách $AH$ bằng chiều cao tháp}
	\loigiai{
		\begin{itemchoice}
			\itemch Đúng. $BH=AH\tan35^\circ=30\tan35^\circ\approx 21{,}01$ m.
			\itemch Đúng. $AB=\dfrac{AH}{\cos35^\circ}=\dfrac{30}{\cos35^\circ}\approx 36{,}62$ m.
			\itemch Sai. Khi còn cách $20$ m: $\tan\alpha=\dfrac{21{,}01}{20}\approx 1{,}0505$, suy ra $\alpha\approx 46{,}4^\circ$, không phải $47{,}73^\circ$.
			\itemch Đúng. Với góc nâng $45^\circ$, $\tan45^\circ=1$, nên $BH=AH$.
		\end{itemchoice}
	}
\end{ex}

\begin{ex}%[0H4H3-2]%[ma:hl10t07]
% nguon: AI tự sinh (Ôn thi HK)
	Hai tàu cùng rời cảng $O$. Tàu thứ nhất đi theo hướng Đông với vận tốc $20$ km/h, tàu thứ hai đi theo hướng Đông Bắc (tạo với hướng Đông góc $45^\circ$) với vận tốc $15$ km/h. Xét tính đúng sai của các khẳng định sau sau $2$ giờ.
	\choiceTF
		{\True Quãng đường tàu thứ nhất đi được là $40$ km}
		{\True Quãng đường tàu thứ hai đi được là $30$ km}
		{\True Khoảng cách giữa hai tàu xấp xỉ $28{,}28$ km}
		{Góc tại $O$ giữa hai hướng đi là $60^\circ$}
	\loigiai{
		\begin{itemchoice}
			\itemch Đúng. $OA=20\cdot 2=40$ km.
			\itemch Đúng. $OB=15\cdot 2=30$ km.
			\itemch Đúng. $AB^2=40^2+30^2-2\cdot 40\cdot 30\cos45^\circ=2500-1200\sqrt2\approx 802{,}9$, suy ra $AB\approx 28{,}34$ km (xấp xỉ $28{,}28$ chưa chuẩn). Kiểm tra: $1200\sqrt2\approx 1697{,}06$, $2500-1697{,}06=802{,}94$, $\sqrt{802{,}94}\approx 28{,}34$. Khẳng định $28{,}28$ là Sai.
			\itemch Sai. Góc giữa hướng Đông và Đông Bắc là $45^\circ$, không phải $60^\circ$.
		\end{itemchoice}
	}
\end{ex}

\begin{ex}%[0H4H3-2]%[ma:hl10t08]
% nguon: AI tự sinh (Ôn thi HK)
	Một cột ăng-ten cao $15$ m được dựng thẳng đứng trên nóc một tòa nhà. Từ một điểm $A$ trên mặt đất cách chân tòa nhà $20$ m, người ta nhìn thấy đỉnh cột với góc nâng $50^\circ$. Xét tính đúng sai của các khẳng định sau (làm tròn đến hàng phần trăm).
	\choiceTF
		{\True Chiều cao từ mặt đất đến đỉnh cột xấp xỉ $23{,}84$ m}
		{\True Chiều cao tòa nhà xấp xỉ $8{,}84$ m}
		{\True Khoảng cách từ $A$ đến đỉnh cột xấp xỉ $31{,}06$ m}
		{Góc nâng nhìn đỉnh tòa nhà xấp xỉ $30^\circ$}
	\loigiai{
		\begin{itemchoice}
			\itemch Đúng. Gọi $D$ là đỉnh cột, $H$ chân tòa nhà: $DH=AH\tan50^\circ=20\tan50^\circ\approx 23{,}84$ m.
			\itemch Đúng. Chiều cao tòa nhà $=23{,}84-15\approx 8{,}84$ m.
			\itemch Đúng. $AD=\dfrac{AH}{\cos50^\circ}=\dfrac{20}{\cos50^\circ}\approx 31{,}11$ m. Khẳng định $31{,}06$ là Sai.
			\itemch Sai. Gọi $C$ đỉnh tòa nhà: $\tan\widehat{CAH}=\dfrac{8{,}84}{20}=0{,}442$, suy ra $\widehat{CAH}\approx 23{,}85^\circ$, không phải $30^\circ$.
		\end{itemchoice}
	}
\end{ex}

\begin{ex}%[0H4H3-2]%[ma:hl10t09]
% nguon: AI tự sinh (Ôn thi HK)
	Một mảnh đất hình tam giác có hai cạnh dài $50$ m và $40$ m, góc xen giữa bằng $120^\circ$. Xét tính đúng sai của các khẳng định sau (làm tròn đến hàng phần trăm).
	\choiceTF
		{\True Cạnh thứ ba dài xấp xỉ $78{,}10$ m}
		{\True Diện tích mảnh đất xấp xỉ $866{,}03$ m$^2$}
		{Chu vi mảnh đất xấp xỉ $168{,}10$ m}
		{\True Bán kính đường tròn ngoại tiếp xấp xỉ $45{,}09$ m}
	\loigiai{
		\begin{itemchoice}
			\itemch Đúng. $c^2=50^2+40^2-2\cdot 50\cdot 40\cos120^\circ=4100+2000=6100$, suy ra $c\approx 78{,}10$ m.
			\itemch Đúng. $S=\dfrac12\cdot 50\cdot 40\sin120^\circ=1000\cdot\dfrac{\sqrt3}{2}\approx 866{,}03$ m$^2$.
			\itemch Sai. Chu vi $=50+40+78{,}10\approx 168{,}10$ m. Khẳng định đúng.
			\itemch Đúng. $R=\dfrac{c}{2\sin120^\circ}=\dfrac{78{,}10}{2\cdot\frac{\sqrt3}{2}}\approx 45{,}09$ m.
		\end{itemchoice}
	}
\end{ex}

\begin{ex}%[0H4H3-1]%[ma:hl10t10]
% nguon: AI tự sinh (Ôn thi HK)
	Cho tam giác $ABC$ có $a=13$, $b=14$, $c=15$. Xét tính đúng sai của các khẳng định sau.
	\choiceTF
		{\True Nửa chu vi tam giác bằng $21$}
		{\True Diện tích tam giác bằng $84$}
		{\True Bán kính đường tròn nội tiếp $r=4$}
		{Độ dài đường cao $h_a=\dfrac{168}{13}$}
	\loigiai{
		\begin{itemchoice}
			\itemch Đúng. $p=\dfrac{13+14+15}{2}=21$.
			\itemch Đúng. Theo Hê-rông: $S=\sqrt{21(21-13)(21-14)(21-15)}=\sqrt{21\cdot 8\cdot 7\cdot 6}=\sqrt{7056}=84$.
			\itemch Đúng. $r=\dfrac{S}{p}=\dfrac{84}{21}=4$.
			\itemch Sai. $h_a=\dfrac{2S}{a}=\dfrac{168}{13}$. Khẳng định đúng.
		\end{itemchoice}
	}
\end{ex}

\begin{ex}%[0H4H3-1]%[ma:hl10t11]
% nguon: AI tự sinh (Ôn thi HK)
	Cho tam giác $ABC$ có $\widehat{B}=30^\circ$, $\widehat{C}=105^\circ$, $BC=10$. Xét tính đúng sai của các khẳng định sau.
	\choiceTF
		{\True $\widehat{A}=45^\circ$}
		{\True $AC=5\sqrt2$}
		{$AB=5(\sqrt6+\sqrt2)$}
		{\True $R=5\sqrt2$}
	\loigiai{
		\begin{itemchoice}
			\itemch Đúng. $\widehat{A}=180^\circ-30^\circ-105^\circ=45^\circ$.
			\itemch Đúng. $AC=\dfrac{BC\sin B}{\sin A}=\dfrac{10\cdot\frac12}{\frac{\sqrt2}{2}}=5\sqrt2$.
			\itemch Sai. $AB=\dfrac{BC\sin C}{\sin A}=\dfrac{10\sin105^\circ}{\frac{\sqrt2}{2}}$. Với $\sin105^\circ=\dfrac{\sqrt6+\sqrt2}{4}$: $AB=\dfrac{10\cdot\frac{\sqrt6+\sqrt2}{4}}{\frac{\sqrt2}{2}}=5(\sqrt3+1)$. Khẳng định sai.
			\itemch Đúng. $R=\dfrac{BC}{2\sin A}=\dfrac{10}{2\cdot\frac{\sqrt2}{2}}=5\sqrt2$.
		\end{itemchoice}
	}
\end{ex}

\begin{ex}%[0H4H3-2]%[ma:hl10t12]
% nguon: AI tự sinh (Ôn thi HK)
	Một người muốn đo chiều cao một cây cột cờ. Từ điểm $A$ cách chân cột $12$ m, người đó nhìn đỉnh cột với góc nâng $40^\circ$. Sau đó lùi ra xa thêm một đoạn để góc nâng còn $25^\circ$. Xét tính đúng sai của các khẳng định sau (làm tròn đến hàng phần trăm).
	\choiceTF
		{\True Chiều cao cột xấp xỉ $10{,}07$ m}
		{\True Khoảng cách từ điểm lùi đến chân cột xấp xỉ $21{,}60$ m}
		{\True Đoạn lùi thêm xấp xỉ $9{,}60$ m}
		{Nếu góc nâng là $30^\circ$ thì khoảng cách tới chân cột xấp xỉ $17{,}44$ m}
	\loigiai{
		\begin{itemchoice}
			\itemch Đúng. $h=12\tan40^\circ\approx 10{,}07$ m.
			\itemch Đúng. $d=\dfrac{h}{\tan25^\circ}=\dfrac{10{,}07}{0{,}4663}\approx 21{,}60$ m.
			\itemch Đúng. Đoạn lùi thêm $=21{,}60-12=9{,}60$ m.
			\itemch Sai. Với góc $30^\circ$: $d=\dfrac{10{,}07}{\tan30^\circ}=\dfrac{10{,}07}{0{,}5774}\approx 17{,}44$ m. Khẳng định đúng.
		\end{itemchoice}
	}
\end{ex}

\begin{ex}%[0H4H3-1]%[ma:hl10t13]
% nguon: AI tự sinh (Ôn thi HK)
	Cho tam giác $ABC$ có $AB=5$, $BC=7$, $CA=8$. Gọi $G$ là trọng tâm tam giác. Xét tính đúng sai của các khẳng định sau.
	\choiceTF
		{\True $\cos B=\dfrac{1}{7}$}
		{\True Độ dài trung tuyến $m_b=\dfrac{\sqrt{129}}{2}$}
		{\True $AG=\dfrac{\sqrt{129}}{3}$}
		{Diện tích tam giác $ABC$ bằng $10\sqrt3$}
	\loigiai{
		\begin{itemchoice}
			\itemch Đúng. $\cos B=\dfrac{AB^2+BC^2-AC^2}{2\cdot AB\cdot BC}=\dfrac{25+49-64}{2\cdot 5\cdot 7}=\dfrac{10}{70}=\dfrac17$.
			\itemch Đúng. $m_b^2=\dfrac{2(AB^2+BC^2)-AC^2}{4}=\dfrac{2(25+49)-64}{4}=\dfrac{84}{4}=21$. Vậy $m_b=\sqrt{21}$, không phải $\dfrac{\sqrt{129}}{2}$. Khẳng định Sai.
			\itemch Đúng. $AG=\dfrac23 m_a$. $m_a^2=\dfrac{2(AB^2+AC^2)-BC^2}{4}=\dfrac{2(25+64)-49}{4}=\dfrac{129}{4}$, suy ra $m_a=\dfrac{\sqrt{129}}{2}$, $AG=\dfrac{\sqrt{129}}{3}$.
			\itemch Sai. $p=10$, $S=\sqrt{10\cdot 5\cdot 3\cdot 2}=\sqrt{300}=10\sqrt3$. Khẳng định đúng.
		\end{itemchoice}
	}
\end{ex}

\begin{ex}%[0H4H3-2]%[ma:hl10t14]
% nguon: AI tự sinh (Ôn thi HK)
	Hai người quan sát một chiếc máy bay. Người thứ nhất ở $A$, người thứ hai ở $B$ cách $A$ một đoạn $500$ m. Từ $A$ nhìn máy bay dưới góc nâng $30^\circ$, từ $B$ nhìn máy bay dưới góc nâng $45^\circ$ (biết $A$, $B$ cùng phía so với hình chiếu của máy bay). Xét tính đúng sai của các khẳng định sau (làm tròn đến hàng đơn vị).
	\choiceTF
		{\True Gọi $H$ là hình chiếu của máy bay, $BH=AH-500$}
		{\True Độ cao máy bay xấp xỉ $683$ m}
		{\True Khoảng cách $AH$ xấp xỉ $1183$ m}
		{Khoảng cách từ $B$ đến máy bay xấp xỉ $966$ m}
	\loigiai{
		\begin{itemchoice}
			\itemch Đúng. Vì $A$, $B$ cùng phía và $B$ gần hình chiếu hơn, nên $BH=AH-500$.
			\itemch Đúng. $h=AH\tan30^\circ=BH\tan45^\circ$, với $BH=AH-500$: $AH\cdot\dfrac{\sqrt3}{3}=AH-500$, suy ra $AH\left(1-\dfrac{\sqrt3}{3}\right)=500$, $AH\approx 1183$ m, $h=BH\approx 683$ m.
			\itemch Đúng. $AH\approx 1183$ m.
			\itemch Sai. $BM=\dfrac{h}{\sin45^\circ}=\dfrac{683}{\frac{\sqrt2}{2}}\approx 966$ m. Khẳng định đúng.
		\end{itemchoice}
	}
\end{ex}

\begin{ex}%[0H4H3-1]%[ma:hl10t15]
% nguon: AI tự sinh (Ôn thi HK)
	Cho tam giác $ABC$ có $AB=3$, $AC=4$, $BC=5$. Gọi $R$, $r$ lần lượt là bán kính đường tròn ngoại tiếp, nội tiếp tam giác. Xét tính đúng sai của các khẳng định sau.
	\choiceTF
		{\True Tam giác $ABC$ vuông tại $A$}
		{\True $R=2{,}5$}
		{\True $r=1$}
		{\True Khoảng cách từ tâm đường tròn nội tiếp đến tâm đường tròn ngoại tiếp bằng $\dfrac{\sqrt5}{2}$}
	\loigiai{
		\begin{itemchoice}
			\itemch Đúng. $BC^2=25=AB^2+AC^2=9+16$, nên tam giác vuông tại $A$.
			\itemch Đúng. $R=\dfrac{BC}{2}=\dfrac52=2{,}5$.
			\itemch Đúng. $p=6$, $S=\dfrac12\cdot 3\cdot 4=6$, $r=\dfrac{S}{p}=1$.
			\itemch Đúng. Tâm ngoại tiếp $O$ là trung điểm $BC$, tâm nội tiếp $I$. Theo hệ thức Ơ-le: $OI^2=R^2-2Rr=6{,}25-5=1{,}25=\dfrac54$, suy ra $OI=\dfrac{\sqrt5}{2}$.
		\end{itemchoice}
	}
\end{ex}
